% META: {"id": "TEXP-ARI-CR-013", "chapitre": "TEXP-ARITHMETIQUE", "type_objet": "cours", "capacites_codes": ["C7", "C8", "C9"], "sources_inspiration": [], "mode_creation": "ex_nihilo", "status": "generated"}

\section{Nombres premiers}

% ============================================================
% STRATE 1 — Infinite des nombres premiers, decomposition
% ============================================================

\definition[1]{
Un entier $p\geq2$ est \textbf{premier} s'il n'admet que deux diviseurs positifs : $1$ et lui-même.
}

\theoreme[Infinité des nombres premiers]{
L'ensemble des nombres premiers est infini.
}

\demonstration{
Par l'absurde, supposons qu'il n'existe qu'un nombre fini de nombres premiers $p_1,p_2,\dots,p_n$. Posons $N=p_1\times p_2\times\dots\times p_n+1$. Cet entier $N$ admet un diviseur premier $p$ (tout entier $\geq2$ admet un diviseur premier). Ce $p$ est nécessairement l'un des $p_i$ (par hypothèse, ce sont les seuls nombres premiers). Mais alors $p_i$ divise à la fois $N$ et $p_1\times\dots\times p_n$, donc $p_i$ divise leur différence $N-p_1\times\dots\times p_n=1$ : c'est impossible, car aucun nombre premier ne divise $1$. Contradiction : l'ensemble des nombres premiers est donc infini.
}

% BEGIN-VERIFY
% premiers = [2, 3, 5, 7]
% N = 1
% for p in premiers:
%     N *= p
% N += 1
% assert N == 211
% # verifie qu'aucun des premiers de la liste ne divise N (illustration de la contradiction)
% for p in premiers:
%     assert N % p != 0
% END-VERIFY

\propriete[Décomposition en facteurs premiers]{
Tout entier $n\geq2$ se décompose de façon \textbf{unique} (à l'ordre près) en produit de nombres premiers.
}

\begin{python}
def facteurs_premiers(n):
    facteurs = []
    d = 2
    while d * d <= n:
        while n % d == 0:
            facteurs.append(d)
            n //= d
        d += 1
    if n > 1:
        facteurs.append(n)
    return facteurs
\end{python}

% BEGIN-VERIFY
% def facteurs_premiers(n):
%     facteurs = []
%     d = 2
%     while d * d <= n:
%         while n % d == 0:
%             facteurs.append(d)
%             n //= d
%         d += 1
%     if n > 1:
%         facteurs.append(n)
%     return facteurs
% assert facteurs_premiers(360) == [2, 2, 2, 3, 3, 5]
% produit = 1
% for f in facteurs_premiers(360):
%     produit *= f
% assert produit == 360
% END-VERIFY

% ============================================================
% STRATE 2 — Petit theoreme de Fermat, crible d'Eratosthene
% ============================================================

\theoreme[Petit théorème de Fermat]{
Si $p$ est premier et $a$ un entier non divisible par $p$, alors $a^{p-1}\equiv1\ [p]$.
}

\exemple{
Pour $p=7$ et $a=3$ : $3^6=729$. On vérifie $729 = 104\times7+1$, donc $3^6\equiv1\ [7]$.
}

% BEGIN-VERIFY
% assert 3**6 == 729
% assert 729 % 7 == 1
% assert pow(3, 6, 7) == 1
% END-VERIFY

\propriete[Crible d'Ératosthène]{
Pour trouver tous les nombres premiers jusqu'à $N$, on part de la liste $2,3,\dots,N$, et on élimine successivement tous les multiples de chaque nombre non encore éliminé, en partant de $2$.
}

\begin{python}
def crible(N):
    est_premier = [True] * (N + 1)
    est_premier[0] = est_premier[1] = False
    for d in range(2, int(N**0.5) + 1):
        if est_premier[d]:
            for multiple in range(d*d, N + 1, d):
                est_premier[multiple] = False
    return [n for n in range(N + 1) if est_premier[n]]
\end{python}

% BEGIN-VERIFY
% def crible(N):
%     est_premier = [True] * (N + 1)
%     est_premier[0] = est_premier[1] = False
%     for d in range(2, int(N**0.5) + 1):
%         if est_premier[d]:
%             for multiple in range(d*d, N + 1, d):
%                 est_premier[multiple] = False
%     return [n for n in range(N + 1) if est_premier[n]]
% assert crible(30) == [2, 3, 5, 7, 11, 13, 17, 19, 23, 29]
% END-VERIFY

\erreurFrequente{
\textbf{Appliquer le petit théorème de Fermat sans vérifier que $p$ est premier.} Le résultat est \textbf{faux} en général si $p$ n'est pas premier : par exemple, pour $p=8$ (non premier) et $a=3$, $3^7=2187=273\times8+3$, donc $3^7\equiv3\ [8]$ et non $1$.
}

% BEGIN-VERIFY
% assert 3**7 == 2187
% assert 2187 % 8 == 3
% END-VERIFY
